\section{Conclusion}\label{sec:concl}

We presented \sname, a lifelong fleet manager for grid warehouses that runs
task allocation, windowed prioritized path planning and charging in one
rolling-horizon loop, couples allocation to planning through arrival times
at the pickup from the same planner, and keeps every epoch conflict-free by
a hold cascade, a variant of existing fail policies, whose conflict-freeness
we proved. An independent trajectory auditor
found no executed conflict in 4{,}740 seeded runs with 60.4 million
robot-steps across four layouts, fleets of 4 to 48 robots, two load
regimes and both planner backends; removing the cascade brings conflicts
back and removing the dynamic priorities gridlocks the fleet. In these runs, with normal
batteries no robot depleted and every run delivered all of its tasks; with
doubled drain the threshold charging policy stranded robots whenever the
slots saturated, which violates the bounded slot wait that our conditional
reserve bound assumes. With dock capacity limiting the tasks in flight, the
pickup-leg coupling did not measurably improve throughput or service time
over static Hungarian or even greedy assignment, neither in the twelve cells
of the main study nor in the regimes in which we expected it to matter that stay inside
the safe envelope. The test is weak: recorded counters show that the
coupled cost changed the matching in only 1--14\% of the allocations of the
main study; only under hotspot demand did it reach 27\%. The study detects no large effect, is not an
equivalence test, and leaves full-trip and dock-aware costs untested.
Dock capacity and the planning window changed throughput
by 7--29\%, and on the small layout a windowed CBS backend shortened the makespan
by 4--12\%. Token passing and token passing with task swaps delivered 16--66\%
less at our default dock capacity, but at
the one robot in flight per dock that their endpoint rule allows, token
passing with task swaps and an interruptible drive home delivered 0.9--3.7\%
more than our loop in every setting, and with the published idle rule,
adapted, it ranged from 10.7\% below to 2.9\% above across the five settings.
The implementation, the benchmark and every number in this paper are
available in the operator console repository so that the study can be
repeated on layouts where the balance may differ.
